\documentclass{article}


% arxiv style file, see https://github.com/kourgeorge/arxiv-style
\usepackage{arxiv}

\usepackage[utf8]{inputenc} % allow utf-8 input
\usepackage[T1]{fontenc}    % use 8-bit T1 fonts
\usepackage{hyperref}       % hyperlinks
\usepackage{url}            % simple URL typesetting
\usepackage{booktabs}       % professional-quality tables
\usepackage{amsfonts}       % blackboard math symbols
\usepackage{nicefrac}       % compact symbols for 1/2, etc.
\usepackage{microtype}      % microtypography
\usepackage{lipsum}		% Can be removed after putting your text content
\usepackage{graphicx}
\usepackage{natbib}
\usepackage{doi}

\usepackage{mathtools}

\title{A template for the \emph{arxiv} style}

%\date{September 9, 1985}	% Here you can change the date presented in the paper title
%\date{} 					% Or removing it

\author{ \href{https://orcid.org/0000-0003-4058-1543}{\includegraphics[scale=0.06]{orcid.pdf}\hspace{1mm}Malte~Schierholz}\thanks{Use footnote for providing further
		information about author (webpage, alternative
		address)---\emph{not} for acknowledging funding agencies.} \\
	Department of Computer Science\\
	Cranberry-Lemon University\\
	Pittsburgh, PA 15213 \\
	\texttt{hippo@cs.cranberry-lemon.edu} \\
	%% examples of more authors
	%% \AND
	%% Coauthor \\
	%% Affiliation \\
	%% Address \\
	%% \texttt{email} \\
	%% \And
	%% Coauthor \\
	%% Affiliation \\
	%% Address \\
	%% \texttt{email} \\
	%% \And
	%% Coauthor \\
	%% Affiliation \\
	%% Address \\
	%% \texttt{email} \\
}

% Uncomment to remove the date
%\date{}

% Uncomment to override  the `A preprint' in the header
%\renewcommand{\headeright}{Technical Report}
%\renewcommand{\undertitle}{Technical Report}
\renewcommand{\shorttitle}{\textit{arXiv} Template}

%%% Add PDF metadata to help others organize their library
%%% Once the PDF is generated, you can check the metadata with
%%% $ pdfinfo template.pdf
\hypersetup{
pdftitle={A template for the arxiv style},
pdfsubject={q-bio.NC, q-bio.QM},
pdfauthor={David S.~Hippocampus, Elias D.~Striatum},
pdfkeywords={First keyword, Second keyword, More},
}

\begin{document}
\maketitle

\begin{abstract}
	\lipsum[1]
\end{abstract}


% keywords can be removed
\keywords{First keyword \and Second keyword \and More}


\section{Introduction}\label{sec:introduction}

Time is a crucial factor in many data sets, but best practices of how date variables should be used in analyses do not exist.

For example, consider the donors-choose data set.\footnote{The KDD Cup 2014 DonorsChoose dataset is available from \href{https://www.kaggle.com/c/kdd-cup-2014-predicting-excitement-at-donors-choose/data}{Kaggle}. The \href{https://github.com/dssg/donors-choose}{Data Science for Social Good Github Repository} lays out a possible route of analysis, our reference method here.} DonorsChoose allows teachers to crowdsource funding for projects, like buying a microwave that lets them cook with children in preschool. The requested items are shipped only when enough citizen donors support the project and funding goals are met within four months. As data scientists, we want to identify postings automatically that are at risk of not meeting its funding goal, and before they get published, with the intention to help teachers improve their postings and funding chances improve. Several projects are posted each day, and it is not unlikely that factors affecting funding success change over time.

A common framework to predict outcomes in settings like DonorsChoose is implemented in triage\footnote{\url{https://pypi.org/project/triage/}}, a python package. The framework is as follows: 
\begin{itemize}
    \item \textit{Training data}: Pool all observations from a given month (or any other user-defined interval), forming the training data. Outcomes are observed up to four months after (again, depending on the task at hand).
    \item \textit{Model training}: Employ any statistical learning technique (e.g., random forest, boosting, linear regression ...) to generate a predictive model.
    \item \textit{Deployment/Test data}: The predictive model gets deployed on future time periods. It is good practice to evaluate it first. For evaluation, all observations from the month after training outcomes became available (or another interval) are pooled, forming the test data.
    \item \textit{Check sensitivity}: Repeat the steps above for different training/test months (or other time intervals), while keeping the time period between both data sets constant.
\end{itemize}
Statistical learning theory provides little guidance here, since a standard theoretical assumption---training and test data are drawn from the same probability distribution---is not met. Instead, the distribution $p_t(y, x)$ might shift over time and in the most extreme case, if there were a sudden shock between training and testing, extrapolation from training to the deployment period would be impossible. The contrary case, that the distribution $p_t(y, x)$ is constant in time, is the implicit working assumption in the above framework, but real data rarely behave this way (as one easily sees with the sensitivity checks). The middle ground is most realistic, namely that $p_t(y, x)$ is continuous in time. While researchers have proposed various ways to correct for distribution shift \citep[e.g.,][p. 133ff.]{varshney_trustworthy_2021}, we are not aware of any research that aims to exploit this assumption of continuity.

\section{Model}\label{sec:model}

Imagine a data set with $i = 1, ..., N$ observations. We want to predict a binary outcome $y$ given predictors $x$. While this situation can be modeled by logistic regression (or any other classifier), our situation differs because we have information about time: All predictors are available at time $t = 1, ..., T$. The outcome will be observed later within a fixed time span (if it were already known, we would not need to predict it). 

Let $y_{i(t)}$ be the outcome of individual $i$, observed within a fixed period after $t$. $x_{i(t)p}$ are the $p = 1, ..., P$ predictors that were available at the time of prediction $t$. (The subscript $t$ indicates that we make predictions at different times, not the time of observation.) $\beta$, $\alpha$ and $\nu$ are parameter matrices, each of dimension $T*P$, with entries for each predictor at all time points.

Our modeling approach is described by the following equations:
\begin{eqnarray}
y_{i(t)} | x_{i(t)p} \sim Bin(1, \textrm{logit}^{-1}(\sum_p x_{i(t)p} \beta_{tp})) & \textrm{(observation equation)} \\
\beta_{tp} \sim N(\alpha_{tp}, \sigma^2_{\beta})\textrm{~~~~~~~~~~~~~~~~~~~~~~~~~~~~~~} & \textrm{(fluctuation equation)} \\
\alpha_{t+1,p} \sim N(\alpha_{t,p} + \nu_{t,p}, \sigma^2_{\alpha})\textrm{~~~~~~~~~~~~~~~~~~} & \textrm{(state equation)} \\
\nu_{t+1,p} \sim N(\nu_{t,p}, \sigma^2_{\eta})\textrm{~~~~~~~~~~~~~~~~~~~~~~~~~~~~~~} & \textrm{(trend equation, optional)}
\end{eqnarray}

The observation equation is a logistic regression model if it were not for the subscript $t$, indicating that parameters $\beta$ are different and depend on the hour/day/year (application-specific) of prediction. There may be very few observations at each time point, impeding the estimation of $\beta$. Since state-space models (also known as dynamic linear models) \citep{prado_time_2010, shumway_time_2011, durbin_time_2012} were developed for situations, in which a single observation (or none) is available at each $t$, they appear well-suited for our analysis. The state equation and the trend equation are well-known concepts in that literature, and many extensions are available from it.

Unlike other approaches we are aware of, we include a fluctuation equation (in state-space terminology this would be called observation equation), only possible because we usually have (far) more than one observation at each time point. The implication is this: If we had at each time point an infinite number of observations, $\sigma^2_\beta$ would go to $\infty$ and the $\beta_t$ would be equal to estimates of $t$ independent logistic regressions. If there are, conversely, just very few observations at each time point, $\sigma^2_\beta$ will turn out to be small and observations from neighboring time points weigh in strongly in the estimation of $\beta_t$.


Warning message with data from 2011:

Warning messages:
1: There were 274 divergent transitions after warmup. See
http://mc-stan.org/misc/warnings.html#divergent-transitions-after-warmup
to find out why this is a problem and how to eliminate them.
2: There were 7726 transitions after warmup that exceeded the maximum treedepth. Increase max_treedepth above 10. See
http://mc-stan.org/misc/warnings.html#maximum-treedepth-exceeded
3: There were 6 chains where the estimated Bayesian Fraction of Missing Information was low. See
http://mc-stan.org/misc/warnings.html#bfmi-low
4: Examine the pairs() plot to diagnose sampling problems
  
5: The largest R-hat is 6.06, indicating chains have not mixed.
Running the chains for more iterations may help. See
http://mc-stan.org/misc/warnings.html#r-hat
6: Bulk Effective Samples Size (ESS) is too low, indicating posterior means and medians may be unreliable.
Running the chains for more iterations may help. See
http://mc-stan.org/misc/warnings.html#bulk-ess
7: Tail Effective Samples Size (ESS) is too low, indicating posterior variances and tail quantiles may be unreliable.
Running the chains for more iterations may help. See
http://mc-stan.org/misc/warnings.html#tail-ess

\section{Data}\label{sec:data}

\section{Analytical Strategy}\label{sec:strategy}

\section{Results}\label{sec:results}

\section{Discussion}\label{sec:discussion}

\section{Latex}
See Section \ref{sec:introduction}.

\subsection{Headings: second level}

\begin{equation}
	\xi _{ij}(t)=P(x_{t}=i,x_{t+1}=j|y,v,w;\theta)= {\frac {\alpha _{i}(t)a^{w_t}_{ij}\beta _{j}(t+1)b^{v_{t+1}}_{j}(y_{t+1})}{\sum _{i=1}^{N} \sum _{j=1}^{N} \alpha _{i}(t)a^{w_t}_{ij}\beta _{j}(t+1)b^{v_{t+1}}_{j}(y_{t+1})}}
\end{equation}


\section{Examples of citations, figures, tables, references}
\label{sec:others}

\subsection{Citations}
Citations use \verb+natbib+. The documentation may be found at
\begin{center}
	\url{http://mirrors.ctan.org/macros/latex/contrib/natbib/natnotes.pdf}
\end{center}

% Here is an example usage of the two main commands (\verb+citet+ and \verb+citep+): Some people thought a thing \citep{kour2014real, hadash2018estimate} but other people thought something else \citep{kour2014fast}. Many people have speculated that if we knew exactly why \citet{kour2014fast} thought this\dots

\subsection{Figures}
\lipsum[10]
See Figure \ref{fig:fig1}. Here is how you add footnotes. \footnote{Sample of the first footnote.}
\lipsum[11]

\begin{figure}
	\centering
	\fbox{\rule[-.5cm]{4cm}{4cm} \rule[-.5cm]{4cm}{0cm}}
	\caption{Sample figure caption.}
	\label{fig:fig1}
\end{figure}

\subsection{Tables}
See awesome Table~\ref{tab:table}.

The documentation for \verb+booktabs+ (`Publication quality tables in LaTeX') is available from:
\begin{center}
	\url{https://www.ctan.org/pkg/booktabs}
\end{center}


\begin{table}
	\caption{Sample table title}
	\centering
	\begin{tabular}{lll}
		\toprule
		\multicolumn{2}{c}{Part}                   \\
		\cmidrule(r){1-2}
		Name     & Description     & Size ($\mu$m) \\
		\midrule
		Dendrite & Input terminal  & $\sim$100     \\
		Axon     & Output terminal & $\sim$10      \\
		Soma     & Cell body       & up to $10^6$  \\
		\bottomrule
	\end{tabular}
	\label{tab:table}
\end{table}

\subsection{Lists}
\begin{itemize}
	\item Lorem ipsum dolor sit amet
	\item consectetur adipiscing elit.
	\item Aliquam dignissim blandit est, in dictum tortor gravida eget. In ac rutrum magna.
\end{itemize}


\bibliographystyle{unsrtnat}
\bibliography{references2}  %%% Uncomment this line and comment out the ``thebibliography'' section below to use the external .bib file (using bibtex) .


%%% Uncomment this section and comment out the \bibliography{references} line above to use inline references.
% \begin{thebibliography}{1}

% 	\bibitem{kour2014real}
% 	George Kour and Raid Saabne.
% 	\newblock Real-time segmentation of on-line handwritten arabic script.
% 	\newblock In {\em Frontiers in Handwriting Recognition (ICFHR), 2014 14th
% 			International Conference on}, pages 417--422. IEEE, 2014.

% 	\bibitem{kour2014fast}
% 	George Kour and Raid Saabne.
% 	\newblock Fast classification of handwritten on-line arabic characters.
% 	\newblock In {\em Soft Computing and Pattern Recognition (SoCPaR), 2014 6th
% 			International Conference of}, pages 312--318. IEEE, 2014.

% 	\bibitem{hadash2018estimate}
% 	Guy Hadash, Einat Kermany, Boaz Carmeli, Ofer Lavi, George Kour, and Alon
% 	Jacovi.
% 	\newblock Estimate and replace: A novel approach to integrating deep neural
% 	networks with existing applications.
% 	\newblock {\em arXiv preprint arXiv:1804.09028}, 2018.

% \end{thebibliography}


\end{document}
